\section{Introduction and Scope}
\label{sec:introduction}

Task 0 selected the code-generation problem: generate a complete Kubernetes
deployment set from one abstract model of what an application needs, rather
than from hand-maintained YAML in which the same fact is written down several
times. This report covers Task 1, the metamodelling step. It describes the
domain, the modelling decisions and the alternatives that were considered, then
presents the two metamodels and the example models that instantiate them.

\subsection{Response to the Task 0 Feedback}

The feedback on Task 0 raised three points, and each changed something in this
report.

\textbf{Scope.} The warning was to be careful with scope, and it was
justified: the original proposal carried three metamodels and an open-ended set
of generated resources. The scope is now fixed by three statements. The pipeline
has \emph{two} metamodels, not three (Section~\ref{sec:decisions}): the
Deliverable Set is generated text, so it is not given a metamodel of its own.
The generated resources are a \emph{closed} set of six generators, listed in
Section~\ref{sec:targetmetamodel}, and adding a seventh is a recorded decision
rather than a new feature. The estate modelled is a \emph{single} cluster, so
every rule that would only become visible at a second cluster is recorded as an
assumption instead of being designed for.

\textbf{Notation.} The question was why UML class diagrams rather than a
deployment or component diagram, and whether ArchiMate or the C4 model would
serve. Section~\ref{sec:notation} answers both: a class diagram is the notation
for a language's abstract syntax and the others describe systems rather than
languages, so they cannot replace it for the metamodels; but for the example
\emph{models} a deployment view is the better picture, and
Section~\ref{sec:examplemodels} now carries one.

\subsection{What Is In Scope and What Is Not}

\begin{table}[!htbp]
\centering
\small
\begin{tabular}{>{\raggedright\arraybackslash}p{0.46\linewidth}>{\raggedright\arraybackslash}p{0.46\linewidth}}
\hline
\textbf{In scope} & \textbf{Out of scope} \\
\hline
Two metamodels: the authored source language, and the resolved target model. &
  A metamodel of the generated files. They are text, and are compared as text. \\
A closed set of six generators for one cluster. &
  Additional target platforms, and a second cluster. \\
Validation of a source model, and the constraints that refuse an invalid one. &
  Applying the generated files to a cluster, which is a separate concern. \\
Dependency resolution between applications. &
  Runtime behaviour, scaling decisions and monitoring rules. \\
\hline
\end{tabular}
\caption{Scope of the Project After the Task 0 Feedback}
\label{tab:scope}
\end{table}

The last exclusion is the one that most affects the report's claims. This
project generates a deployment set; it does not deploy it. Every statement about
correctness in this report is therefore a statement about the generated model
and the generated files, never about a running cluster.
